% !TeX root = ../../tesis.tex

% =============================================================================
%  B.5 — CLIENTE HTTP: PUENTE CON APIMETRO
% =============================================================================
\section{Cliente HTTP: Puente con el Servidor Apimetro}
\label{anx:vft-goclient}

La capa de infraestructura de VFTModel incluye un cliente asíncrono
(\path{src/infrastructure/go_client/client.py}) que consume los endpoints
espaciales del servidor Go de Apimetro. La URL base se configura mediante la
variable de entorno \texttt{APIMETRO\_URL}, definida en
\path{src/infrastructure/go_client/settings.py}, de modo que el mismo código
funciona tanto en desarrollo local como apuntando al servidor remoto.

\begin{lstlisting}[style=estiloPython,
    caption={Configuración de la URL de Apimetro por variable de entorno.
             \texttt{src/infrastructure/go\_client/settings.py}.},
    label=lst:vft-settings]
import os
from dotenv import load_dotenv

load_dotenv(os.getenv("ENV_FILE", ".env.local"))
APIMETRO_URL = os.getenv("APIMETRO_URL", "http://localhost:8080/movilidad")
\end{lstlisting}

\subsection*{Descarga paralela de estaciones y líneas}

La función \texttt{fetch\_full\_network} descarga estaciones y líneas en
concurrencia con \texttt{asyncio.gather} y las unifica en un único
\texttt{FeatureCollection} listo para el validador Pydantic. Si ambas
peticiones fallan ---por ejemplo, cuando Apimetro no está disponible--- el
código activa un protocolo de \termino{fallback} que lee el archivo
\texttt{map.geojson} almacenado localmente.

\begin{lstlisting}[style=estiloPython,
    caption={Descarga paralela y unificación de capas espaciales de Apimetro.
             \texttt{src/infrastructure/go\_client/client.py}.},
    label=lst:vft-goclient]
async def fetch_full_network() -> dict:
    url_estaciones = f"{APIMETRO_URL}/mapas/geojsonEstacion"
    url_lineas     = f"{APIMETRO_URL}/mapas/geojsonLinea"

    resultados = await asyncio.gather(
        fetch_single_geojson(url_estaciones),
        fetch_single_geojson(url_lineas)
    )
    features_unificados = []
    for geojson in resultados:
        payload = geojson.get("data", geojson)
        features_unificados.extend(payload.get("features", []))

    # Protocolo Fallback: si la red Go no responde, usar map.geojson local
    if not features_unificados:
        vft_logger.warning("Activando protocolo Fallback (map.geojson local).")
        with open(os.path.join(os.getcwd(), "map.geojson"), "r") as f:
            features_unificados = json.load(f).get("features", [])

    return {"type": "FeatureCollection", "features": features_unificados}
\end{lstlisting}

\subsection*{Manejo de errores por endpoint}

Cada petición individual (\texttt{fetch\_single\_geojson}) captura dos tipos
de falla: errores de conexión (\texttt{httpx.RequestError}) y errores de
formato \textsc{json}. En ambos casos devuelve un
\texttt{FeatureCollection} vacío en lugar de propagar la excepción,
permitiendo que \texttt{fetch\_full\_network} continúe con los datos
disponibles del otro endpoint.
